
\chapter{Accessing databases}

\index{database}
\index{DBMS}

``Database'' is a very loose term for ``a repository of information that can
be readily accessed.'' When we use that term in Data Science, we mean a data
source that is managed by a piece of software called a \textbf{DBMS}, or
\textbf{database management system}. (True, plain-old lists, dictionaries, or
Pandas \texttt{DataFrame}s are also ``repositories of information that can be
accessed,'' but when we use the term \textbf{database} we're not talking about
that.)

There are principally two reasons to learn how to access data from databases:

\begin{itemize}
\itemsep.1em
\item Sometimes, our data is simply available in that form only. We may want
to import it into Python/Pandas and deal with it using the techniques from the
first half of the semester, but in order to get it in there, we need to
(briefly) interface with the database and issue the proper commands.
\item Sometimes, our data is just flat \textit{big}. There's a limit to how
much data Spyder can store in memory and deal with on the spot. Suppose
you had to analyze all the eBay transactions over the past two years? Or all
the tweets made in the world over the last month? That stuff can't remotely
fit in your system's memory. Hence, we have to access it through a database
which will manage the sheer size and scope for us, and work with it a little
at a time.
\end{itemize}

\section{Two kinds of databases}

\index{relational database}
\index{NoSQL database}

There are essentially two kinds of databases in the world:
\textbf{relational}, and \textbf{NoSQL}.

\index{table!relational database}
\index{SQL}
\index{RDBMS}

Previously this semester, we've worked with data spread across multiple
\texttt{DataFrame}s that we need to join together in some way. This structure
is extremely common, so much so that most of traditional database schema theory
is based directly on it. A \textbf{relational database} is a repository of
information in which the data is stored in \textbf{tables} (a table is pretty
much \textit{exactly} a \texttt{DataFrame} in spirit), and on which the
language \textbf{SQL}\footnote{Some people pronounce the acronym SQL as the
three letters ``ess-queue-ell,'' and others as the word ``sequel.'' Either one
is acceptable (I do the former).} can be used to query it and manipulate it. A
software system that enables a user to create and work with relational
databases is called an \textbf{RDBMS}, or \textbf{relational database
management system}.

\index{Oracle}
\index{Sybase}
\index{Microsoft SQL Server}
\index{MySQL}
\index{MariaDB}
\index{PostgreSQL}
\index{Codd, Ted}

There are many, many RDBMS's out there: Oracle, Sybase, Microsoft SQL Server,
MySQL, MariaDB, PostgreSQL, and literally dozens of others. They all work in
very much the same way: they store data in tables, and they are manipulated via
SQL. Ever since Ted Codd invented the relational model in 1970, this
organizational paradigm has proven to be so powerful and computationally
efficient that it has dominated the database world. Only recently have some
alternate paradigms gained some traction (so-called ``NoSQL'' databases), but
even so, the majority of the world's structured data is all in relational
databases of some sort.

\index{Cassandra}
\index{MongoDB}
\index{Redis}
\index{Neo4j}
\index{key-value pair}
\index{graph (network)}

Relational DBMS's are much more similar to each other than NoSQL systems are to
each other. The wide variety of NoSQL DBMS's (Cassandra, MongoDB, Redis, Neo4j,
\textit{etc.}) means that learning how to use one doesn't give you much
knowledge about how to use one of the others. Some store key/value pairs;
others store objects and binary data; some store information as a graph
(network). Part of the great power of the relational model is that SQL has
proven to be a common, flexible, broadly-applicable query language that many
different vendors have adopted and which suits many data needs.

\section{SQLite}

\index{SQLite}
\index{RDBMS}

SQLite is a really slim, trim RDBMS that stores all of its data \textit{in a
single file.} When you work with SQLite (or any RDBMS), though, you don't think
in terms of files, but in terms of tables. The fact that the tables are stored
in file(s) in some strange format is invisible to you as the user of the data,
and this is in fact one of the great powers of the relational model: separating
the conceptual data from the storage implementation.

\subsection{Interfacing to SQLite}

\index{connect@\texttt{connect()} (SQLite)}

``\textbf{Connecting}'' to a relational database basically means establishing a
communication path over which you can talk SQL. In Python, this is pretty easy:

\begin{Verbatim}[fontsize=\small,samepage=true,frame=single,framesep=3mm]
import sqlite3

conn = sqlite3.connect('my_local_file.db')
\end{Verbatim}
\label{connectSqlite}

\index{extension (filename)}
\index{db file@\texttt{.db} file}

As indicated, SQLite-compatible files typically have a \texttt{.db} extension
(for ``database.'')

\index{query (database)}

The return value from \texttt{connect()} that we're calling \texttt{conn} is
``a connection to the database.'' We'll use it every time we want to perform an
\textbf{SQL query}. SQL is a new language, but it's not nearly as complex as
something like Python because it's not a programming language, but a
\textbf{query language}. You don't write entire programs in it; you merely
specify queries in it. This means that the range of things you can do in SQL is
much more restricted, and hence the language is easier to learn.

\subsection{Browsing the tables}

\index{read\_sql@\texttt{read\_sql()} (Pandas)}
\index{table!relational database}

The first thing you'll normally want to do is figure out which tables exist in
the database. You can run this command:

\index{sqlite\_master@\texttt{sqlite\_master}}
\label{sqlite_master}
\index{SELECT statement@\texttt{SELECT} statement (SQL)}
\index{FROM clause@``\texttt{FROM}'' clause (SQL)}

\begin{Verbatim}[fontsize=\footnotesize,samepage=true,frame=single,framesep=3mm]
print(pd.read_sql(
    "SELECT name FROM sqlite_master WHERE type='table'", conn))
\end{Verbatim}

to find out. If this were the output:

\begin{Verbatim}[fontsize=\small,samepage=true,frame=leftline,framesep=5mm,framerule=1mm]
          name
0        songs
1       albums
2      artists
3       labels
\end{Verbatim}

that would tell us that the SQLite database has four tables in it, named
\texttt{songs}, \texttt{albums}, \texttt{artists}, and \texttt{labels}. Each
one can be queried using SQL statements, or simply imported wholesale into a
Pandas \texttt{DataFrame} (see below).

\section{Querying the database}

\index{read\_sql@\texttt{read\_sql()} (Pandas)}

By far the most common thing you'll do with a relational database like SQLite
is execute an SQL query, and get a \texttt{DataFrame} back:

\begin{Verbatim}[fontsize=\small,samepage=true,frame=single,framesep=3mm]
df_name = pd.read_sql("YOUR SQL QUERY GOES HERE", conn)
\end{Verbatim}

\index{embedded statement@``embedded'' statement}

The string inside the call to \texttt{pd.read\_sql()} is called ``an embedded
SQL statement.'' This means we have code in one language (SQL) ``embedded''
inside another (Python). Do not be alarmed at this.


\subsection{The easy case: inhaling an entire table}
\index{the easy case}
\label{easy1}

Now \textit{if} the data is small enough to load into memory, you don't need to
know much SQL at all. You can simply use this one-liner:

\index{SELECT statement@\texttt{SELECT} statement (SQL)}
\index{*@\texttt{*} (splat) (SQL)}
\index{FROM clause@``\texttt{FROM}'' clause (SQL)}

\begin{Verbatim}[fontsize=\small,samepage=true,frame=single,framesep=3mm]
df = pd.read_sql("SELECT * FROM table_name", conn)
\end{Verbatim}

This reads one of the database's tables (whose names you learned from the
command in Section~\ref{sqlite_master}) into a new, fully-fledged
\texttt{DataFrame}, with column names, a reasonable guess at their types,
\textit{etc.} 

That's literally all you need to know about working with a (smallish) SQLite
database. Just to connect to it, get all the table names, read each one into a
\texttt{DataFrame}, and continue in Pandas Land, forgetting all about SQLite.
You can then leverage all the knowledge you've learned up to this point. Bingo!


\subsection{The harder case: more complex SQL}

\index{SQL}

Sometimes, however, the data is too large to do this. So it's worth learning a
little SQL so you can make the RDBMS do some of the heavy lifting, and give you
trimmed results for your Python program to work with.

\index{SELECT statement@\texttt{SELECT} statement (SQL)}

The workhorse is the SQL \texttt{SELECT} statement. Here are some common
variants:

\index{*@\texttt{*} (splat) (SQL)}
\begin{itemize}
\itemsep1em
\item To compute the \textit{number} of rows, or individual values, without
returning the entire result, use \texttt{COUNT(*)}, \texttt{MAX()},
\texttt{MIN()}, \textit{etc.}:

\index{COUNT@\texttt{COUNT()} (SQL)}
\index{MAX@\texttt{MAX()} (SQL)}
\index{MIN@\texttt{MIN()} (SQL)}
\index{FROM clause@``\texttt{FROM}'' clause (SQL)}
\begin{BVerbatim}[fontsize=\small,samepage=true,frame=single,framesep=3mm,bgcolor=AliceBlue]
SELECT COUNT(*) FROM table_name
SELECT MAX(col1) FROM table_name
\end{BVerbatim}

\index{LIMIT clause@``\texttt{LIMIT}'' clause (SQL)}
\item To get only the first few rows (similar to Pandas' \texttt{.head()}
operation) tack on a ``\texttt{LIMIT}'' expression:

\begin{BVerbatim}[fontsize=\small,samepage=true,frame=single,framesep=3mm,bgcolor=AliceBlue]
SELECT * FROM table_name LIMIT 10
\end{BVerbatim}

\item To get only certain columns, change the ``\texttt{*}'' to a
comma-separated list of column names:

\begin{BVerbatim}[fontsize=\small,samepage=true,frame=single,framesep=3mm,bgcolor=AliceBlue]
SELECT col1 FROM table_name
SELECT col1,col2,col3 FROM table_name LIMIT 10
\end{BVerbatim}

\index{WHERE clause@``\texttt{WHERE}'' clause (SQL)}
\item To get only certain rows, add a ``\texttt{WHERE} clause'' after the table
name:

\begin{BVerbatim}[fontsize=\footnotesize,samepage=true,frame=single,framesep=3mm,bgcolor=AliceBlue]
SELECT * FROM table_name WHERE col1==val1
SELECT col2 FROM table_name WHERE col1==val1 AND col2>val2
SELECT c1,c2 FROM t WHERE c1==val1 AND NOT c2>val2 LIMIT 10
\end{BVerbatim}

(The usual boolean operations work with keywords \texttt{AND}, \texttt{OR},
and \texttt{NOT}, and parentheses for grouping.)

\index{IN keyword@``\texttt{IN}'' keyword (SQL)}
\item The ``\texttt{IN}'' keyword can be useful in a \texttt{WHERE} clause
for querying membership in a set:

\begin{BVerbatim}[fontsize=\small,samepage=true,frame=single,framesep=3mm,bgcolor=AliceBlue]
SELECT * FROM table_name WHERE col1 IN (val1,val2)
SELECT col1,col2 FROM table_name WHERE col1==val1 AND
            NOT col2 IN (val1,val2,val3) LIMIT 10
\end{BVerbatim}

\index{unique@unique rows (with \texttt{DISTINCT})}
\index{DISTINCT@\texttt{DISTINCT} (SQL)}
\item To get only unique rows (like Pandas' \texttt{.drop\_duplicates()}
gives) use the word \texttt{DISTINCT} before the column names:

\begin{BVerbatim}[fontsize=\small,samepage=true,frame=single,framesep=3mm,bgcolor=AliceBlue]
SELECT DISTINCT * FROM table_name
SELECT col1,col2 FROM table_name WHERE NOT col1 IN
            (val1,val2,val3) OR col2!=val4 LIMIT 5
\end{BVerbatim}

\index{DISTINCT@\texttt{DISTINCT} (SQL)}
\index{COUNT@\texttt{COUNT()} (SQL)}
The \texttt{DISTINCT} option is useful in conjunction with \texttt{COUNT()}:

\begin{BVerbatim}[fontsize=\small,samepage=true,frame=single,framesep=3mm,bgcolor=AliceBlue]
SELECT COUNT(DISTINCT col1) FROM table_name
\end{BVerbatim}

\index{ORDER BY@\texttt{ORDER BY} (SQL)}
\item You can have SQLite do sorting with ``\texttt{ORDER BY}'':

\index{DESC@\texttt{DESC} (``descending'') (SQL)}
\begin{BVerbatim}[fontsize=\small,samepage=true,frame=single,framesep=3mm,bgcolor=AliceBlue]
SELECT * FROM table_name ORDER BY col1
SELECT DISTINCT c1,c2 FROM table_name WHERE c1==val1
            ORDER BY c2,c1 DESC LIMIT 5
\end{BVerbatim}

As you can see, you can sort on multiple columns, and sort in reverse order
with the \texttt{DESC} keyword.

\index{GROUP BY@\texttt{GROUP BY} (SQL)}
\index{groupby@\texttt{.groupby()} (Pandas)}
\item SQL supports a ``GROUP BY'' operation very much like Pandas'
\texttt{.groupby()} method:

\begin{BVerbatim}[fontsize=\small,samepage=true,frame=single,framesep=3mm,bgcolor=AliceBlue]
SELECT col1,count(*) FROM table_name GROUP BY col1
SELECT col1,max(col2) FROM table_name WHERE col3==val1
            GROUP BY col2 ORDER BY col1 DESC LIMIT 5
\end{BVerbatim}

\index{join!JOIN@\texttt{JOIN} (SQL)}
\item Finally, to do a merge/join operation, use the \texttt{JOIN} keyword:

\begin{BVerbatim}[fontsize=\small,samepage=true,frame=single,framesep=3mm,bgcolor=AliceBlue]
SELECT * FROM table_1 join table_2 ON col1==col2
SELECT c1,c2 FROM t1 join t2 ON t1.x==t2.q
            WHERE c2>val1 LIMIT 5
\end{BVerbatim}

The \texttt{JOIN} keyword supports \texttt{LEFT}, \texttt{RIGHT}, and
\texttt{OUTER} options.

\end{itemize}

\index{Oracle}
\index{Sybase}
\index{Microsoft SQL Server}
\index{MySQL}
\index{MariaDB}
\index{PostgreSQL}
\index{Informix}

There are many, many other options as well. See the online documentation at
\url{https://www.sqlite.org/docs.html}. Again, one of the great things about
SQL is that with the exception of the initial commands to connect and to list
all table names (p.~\pageref{sqlite_master}), \textit{all relational databases
use essentially identical syntax.} So if you ever need to query an Oracle,
Sybase, Informix, MySQL, PostgreSQL, or SQL Server database, you'll discover
that everything in this section applies nearly identically.

